This section describes the datasets used in this study, the DE-LSTM baseline
against which the hybrid is developed, and the evaluation protocol shared by all
experiments.

\subsection{Datasets}\label{sec:datasets}

\subsubsection{SIT LiFePO$_4$ Dataset (Primary)}
The SIT dataset~\cite{ponnambalam2026sitdata} comprises 20 commercial
LiFePO$_4$ prismatic cells (Narada Power Source, 50\,Ah nominal, 3.2\,V)
cycled in the Singapore Institute of Technology Battery Laboratory over
18 months.
Cells are organised into three groups differing in temperature and C-rate:
\begin{itemize}
  \item \textbf{G1} (10 cells, IDs 001-1 to 001-8, 101-1, 101-3):
        ambient temperature (${\approx}25$\,\textcelsius), 1\,C charge/discharge,
        deep cycling. All cells share nominally identical specifications
        and protocol, so observed differences are most plausibly dominated by
        manufacturing tolerances, although cycler-channel differences, contact
        resistance, minor temperature variation, testing interruptions and
        production batch may also contribute.
  \item \textbf{G2} (4 cells, IDs 002-1 to 002-4):
        40\,\textcelsius\ thermal chamber, 1\,C rate.
  \item \textbf{G3} (6 cells, IDs 002-5, 002-7, 003-1, 003-3, 003-5, 003-7):
        40\,\textcelsius\ thermal chamber, 2\,C rate.
\end{itemize}
SOH is computed from discharge capacity as
$\text{SOH}(k) = Q_k / Q_1$,
where $Q_k$ is discharge capacity at cycle $k$ and $Q_1$ is the first-cycle
capacity.
Table~\ref{tab:sit_cells} summarises cycle counts and final SOH for all
20 cells, where SOH$_\text{end}$ is the final-cycle SOH.
Cell 003-5 records SOH$>$1, a capacity-recovery artefact arising at low C-rate
after elevated-temperature cycling; it is retained in the analysis but treated
with caution.

\begin{table}[pos=htbp]
\caption{SIT cell inventory across three operating conditions.}
\label{tab:sit_cells}
\resizebox{\columnwidth}{!}{%
\begin{tabular}{llrrr}
\toprule
Group & Cell & Cycles & SOH$_\text{end}$ & Condition \\
\midrule
G1 & 001-1 & 876  & 0.767 & \multirow{10}{*}{25\,\textcelsius, 1\,C} \\
   & 001-2 & 764  & 0.878 & \\
   & 001-3 & 788  & 0.493 & \\
   & 001-4 & 878  & 0.709 & \\
   & 001-5 & 684  & 0.885 & \\
   & 001-6 & 876  & 0.854 & \\
   & 001-7 & 876  & 0.885 & \\
   & 001-8 & 764  & 0.067 & \\
   & 101-1 & 864  & 0.696 & \\
   & 101-3 & 697  & 0.118 & \\
\midrule
G2 & 002-1 & 564  & 0.924 & \multirow{4}{*}{40\,\textcelsius, 1\,C} \\
   & 002-2 & 572  & 0.888 & \\
   & 002-3 & 561  & 0.915 & \\
   & 002-4 & 560  & 0.660 & \\
\midrule
G3 & 002-5 & 689  & 0.913 & \multirow{6}{*}{40\,\textcelsius, 2\,C} \\
   & 002-7 & 695  & 0.564 & \\
   & 003-1 & 603  & 0.461 & \\
   & 003-3 & 693  & 0.559 & \\
   & 003-5 & 567  & 1.075$^\dagger$ & \\
   & 003-7 & 694  & 0.587 & \\
\bottomrule
\multicolumn{5}{l}{$\dagger$ Capacity recovery artefact; included in analysis with caution.}
\end{tabular}%
}
\end{table}

G1 provides the cell-to-cell variability experiment (Section~\ref{sec:variability});
G1/G2/G3 provide the cross-condition transfer experiment
(Section~\ref{sec:brittleness}).

\subsubsection{MIT LFP Dataset (Pretraining)}
The Severson et al.~\cite{severson2019data} dataset provides 124 LFP/graphite
18650 cells (1.1\,Ah) cycled at Stanford University and Toyota Research
Institute at various C-rates and charging protocols.
Of these, 49 cells have complete cycle-level capacity data in our processing;
the \emph{pretraining corpus} comprises the 40 batch-1 cells
(batch1\_cell000 to 039).
This choice is empirical: in a corpus-sensitivity test, extending the corpus
with the nine remaining complete-data cells (including the two batch-3 cells
with $\leq$9\% total fade) collapsed the next-cycle pretraining objective
toward a trivial persistence solution (validation MSE
$3.4\times10^{-3}\rightarrow1\times10^{-5}$) and degraded downstream transfer;
the 40-cell corpus is therefore retained, and the full corpus-sensitivity
results are provided in the released repository.
We emphasise that this is a property of the pretraining objective, not a
fragility of the deployed model: the hybrid is evaluated on, and performs well
on, slowly fading target cells (e.g.\ the G2 cells with final SOH
${\approx}0.9$).
The early-life correlation analysis of Table~\ref{tab:fingerprint} uses the
47 genuine degraders (final SOH $<1$) among all 49 complete-data cells,
excluding the two capacity-recovery artefacts.
Pretraining on chemistry matched LFP cells is a deliberate design choice
motivated by the chemistry and domain alignment effect evaluated in
Section~\ref{sec:discussion}.

\subsubsection{External Datasets}
Four further datasets support the cross-chemistry test and the transfer and RUL
experiments. The NASA Prognostics Center of Excellence
dataset~\cite{saha2007battery} provides 26 lithium cobalt oxide (LCO) 18650
cells of about 2\,Ah cycled at several ambient temperatures under
constant-current discharge; because its chemistry differs from the LFP target
cells, it serves as the cross-chemistry generalisation test. The LISHEN
dataset~\cite{lishen_zenodo} provides 3 large-format 40\,Ah LFP prismatic cells
cycled at 2\,C over a 20\% depth-of-discharge window, used as a third,
independent LFP source in the cross-dataset transfer matrix. The HUST~\cite{ma2022real} and Sandia
(SNL)~\cite{preger2020degradation} collections provide 77 and 29 LFP 18650 cells respectively; together with
the 49 MIT cells they form the 155-cell LFP foundation corpus that pretrains the
RUL decoder (Section~\ref{sec:rul_foundation}). Table~\ref{tab:datasets}
consolidates the chemistry, format, capacity, cell count, operating conditions,
and role of every dataset used in this study.

\begin{table*}[pos=tbp]
\centering
\caption{Datasets used in this study.}\label{tab:datasets}
\footnotesize
\resizebox{\textwidth}{!}{%
\begin{tabular}{l l l r r l l}
\toprule
Dataset & Chemistry & Format & Cap.\ (Ah) & Cells & Conditions & Role in this study \\
\midrule
SIT (primary)~\cite{ponnambalam2026sitdata} & LFP & Prismatic & 50  & 20 & 25 to 40\,\textcelsius, 1 to 2\,C     & Target cells: testing and checkpoint selection \\
MIT/Severson~\cite{severson2019data}  & LFP & 18650     & 1.1 & 40 & Fast charge, multi-protocol           & Base-hybrid pretraining (training) \\
NASA PCoE~\cite{saha2007battery}     & LCO & 18650     & 2   & 26 & Varied temperature, constant current  & Cross-chemistry test (testing) \\
LISHEN~\cite{lishen_zenodo}        & LFP & Prismatic & 40  & 3  & 2\,C, 20\% depth of discharge         & Third LFP source, cross-dataset transfer \\
HUST~\cite{ma2022real}          & LFP & 18650     & 1.1 & 77 & Multi-rate discharge                  & RUL foundation corpus (training) \\
SNL (Sandia)~\cite{preger2020degradation}  & LFP & 18650     & 1.1 & 29 & Varied temperature, rate, DOD         & RUL foundation corpus (training) \\
\bottomrule
\end{tabular}%
}
\end{table*}

\subsection{Baseline: DE-LSTM}\label{sec:delstm}

DE-LSTM~\cite{ponnambalam2026delstm} is the strongest model from our prior work
and serves as the baseline against which the hybrid is developed.
It is a plain LSTM network whose
hyperparameters (hidden units $H$, look-back window $W$,
learning rate $\eta$) are optimised by differential
evolution (DE)~\cite{storn1997differential} using
2-fold walk-forward cross-validation as the fitness function;
the batch size is set deterministically to $\max(32, H)$ rather than searched.
The architecture is:
\begin{equation}
  \hat{y}_t = \mathbf{W}_o\, h_t + b_o,\quad
  h_t, c_t = \mathrm{LSTM}(x_{t-W:t},\, h_{t-1}, c_{t-1})
\end{equation}
where $x_{t-W:t}$ is the SOH look-back window of length $W$, $h_t$ and $c_t$ are
the LSTM hidden and cell states at cycle $t$, $\mathbf{W}_o$ and $b_o$ are the
output-layer weight matrix and bias, and $\hat{y}_t$ is the predicted next-cycle
SOH.

For the SIT baseline evaluation (Section~\ref{sec:baseline}), DE is run
independently for each cell and each training split, with the search space
$H \in \{8, 16, 32, 64, 128, 256\}$, $W \in [3, 30]$, $\eta \in [10^{-4}, 0.05]$.
The DE search uses a population of 3 per parameter and a maximum of
8 generations (\texttt{polish=False}, \texttt{tol}~$=10^{-4}$), on the order of
80 candidate evaluations per cell and split, and is itself run under a single
fixed seed.
The selected hyperparameters are then re-evaluated under two independent random
seeds, and the mean RMSE is reported.

\subsection{Evaluation Protocol}\label{sec:eval}

All experiments use three training splits: 30\%, 50\%, and 70\% of each
cell's SOH trajectory, with the remainder held out for evaluation.
Because each split is defined as a fraction of a cell's already-completed
trajectory, these splits are retrospective (they presuppose the cell's total
lifetime). Prediction is performed in a rolling one-step-ahead manner: at each
held-out cycle the model receives the \emph{observed} previous $W$ SOH values and
predicts the next one, and predicted values are never fed back as inputs (i.e.,
we report one-step-ahead prediction, not recursive multi-step forecasting).
Hybrid-model evaluations are averaged over five random seeds per cell;
single-cell baseline methods (Simple, PA-, GA-, DE-LSTM, TCN, Transformer)
are averaged over two seeds after their hyperparameter search, which itself
runs under a single fixed seed.
Summary tables report the mean RMSE with its
standard deviation across the cells in each group.
Improvement over scratch LSTM is computed per cell and then averaged, so a
reported improvement is the mean of the per-cell percentage RMSE reductions rather
than the percentage change between the group-mean RMSEs.

\subsubsection{Data Partitioning and Leakage}
Within each target cell the training and evaluation partitions are contiguous in
cycle order: training windows are drawn from the first $p\%$ of the trajectory and
evaluation windows from the remainder, so no held-out SOH value is ever used as a
training target. The first evaluation window is seeded with the last $W$
\emph{observed} training cycles as its input; this is inherent to one-step-ahead
prediction and exposes no held-out target. For pretraining, the look-back windows
of all corpus cells are pooled in cell order and the final 10\% of the pooled array
is held out for early stopping. Because the pooled array is not shuffled before
this split, the validation slice comprises whole cells except for the single cell
straddling the boundary, which is consequently the only cell able to contribute
(mutually overlapping) windows to both partitions. Finally, since a cell of length
$T$ yields $T-W$ windows, cells with longer trajectories contribute proportionally
more pretraining samples: the pooled corpus is sample-weighted by trajectory
length rather than balanced per cell.

\subsubsection{Calibration Cells and Test Cells}
Three SIT cells, the numerically first cell of each condition group
(001-1, 002-1, 002-5), designated \emph{a priori}, are reserved as
\emph{calibration cells} for pretraining-checkpoint selection
(Section~\ref{sec:selection}).
All reported SIT results are computed on the remaining \textbf{17 test cells};
the calibration cells never contribute to any reported metric.

\subsubsection{Reproducibility}
All experiments run with deterministic GPU kernels
(TensorFlow op-determinism, full float32) and complete seeding of the Python,
NumPy, and TensorFlow generators, so every reported number is bit-for-bit
reproducible on the same hardware and software stack.
The code, the full pipeline, and the selected pretrained checkpoint are
permanently archived on Zenodo~\cite{ponnambalam2026code}.
Statistical significance of pairwise comparisons uses the Wilcoxon
signed-rank test~\cite{wilcoxon1945individual} at $\alpha = 0.05$.
This non-parametric test is preferred over a paired $t$-test because
per-cell RMSE distributions are skewed by outlier cells (e.g., catastrophic
degraders such as cell 001-8), and the Wilcoxon test does not assume normality.
For the hybrid model, empirical 95\% prediction-interval (PI) coverage (fraction
of test points within the predicted interval) and mean PI width are additionally
reported. These MC-Dropout intervals are approximate prediction intervals, not
calibrated prediction intervals in the conformal sense.

\subsubsection{Model Configurations and Tuning Budget}
Table~\ref{tab:hyperparams} reports every model's architecture, tuning budget, and
training configuration exactly as run. The three search-based LSTMs (PA-, GA- and
DE-LSTM) each receive a per-cell, per-split hyperparameter search of comparable
size (of the order of 80 candidate evaluations), whereas the Simple LSTM, the TCN
and the Transformer share a common fixed configuration. We note this asymmetry
explicitly: the non-recurrent baselines are not hyperparameter-optimised, so their
weaker performance should be read as evidence that a richer backbone is not what
these data reward rather than as a tuned upper bound. The proposed hybrid likewise
receives \emph{no} per-cell search, and so competes against the searched LSTM
baselines from a single fixed configuration.

\begin{table*}[pos=tbp]
\caption{Model configurations and training budgets, as run (per cell and split).}
\label{tab:hyperparams}
\centering
\footnotesize
\begin{tabularx}{\textwidth}{l X l X}
\toprule
Model & Architecture & Tuning budget & Training configuration \\
\midrule
Simple LSTM & LSTM(64) $\to$ Dense(1) & none (fixed) & $W{=}20$, $\eta{=}10^{-2}$, batch 64, 200 epochs, early stop patience 10 \\
PA-LSTM & LSTM $\to$ attention $\to$ pooling $\to$ Dense(1) & PSO, ${\approx}80$ evals & searched: $H\in\{8\ldots256\}$, $W\in[3,30]$, $\eta\in[10^{-4},5{\times}10^{-2}]$; batch $\max(32,H)$, 200 epochs, patience 10 \\
GA-LSTM & LSTM $\to$ Dense(1) & GA, $9{\times}9$ evals & searched, as PA-LSTM \\
DE-LSTM & LSTM $\to$ Dense(1) & DE, ${\approx}80$ evals & searched, as PA-LSTM \\
TCN & 3 residual dilated causal conv blocks (64 filters, $k{=}3$, dilations $1,2,4$) & none (fixed) & $W{=}20$, $\eta{=}10^{-2}$, batch 64, 200 epochs, patience 10 \\
Transformer & 1 encoder block ($d_{\mathrm{model}}{=}64$, 4 heads, feed-forward 128) & none (fixed) & $W{=}20$, $\eta{=}10^{-3}$, batch 64, 200 epochs, patience 10 \\
\midrule
Scratch LSTM & LSTM(64) $\to$ Dropout(0.2) $\to$ Dense(1) & none (fixed) & $W{=}20$, $\eta{=}10^{-3}$, 100 epochs, patience 10 \\
Hybrid (ours) & CNN fingerprint (Conv1D 32 and 16, $k{=}3$, global average pooling, Dense 16) $\to$ LSTM(64) initialised with $(h_0,c_0)$ $\to$ MC-Dropout(0.2) $\to$ Dense(1); 21\,089 parameters & none (fixed) & pretrain: $\eta{=}10^{-3}$, batch 64, 100 epochs, patience 12; fine tune: $\eta{=}5{\times}10^{-4}$, 50 epochs, patience 8; 50 MC-Dropout passes \\
\bottomrule
\end{tabularx}
\end{table*}
